\reading{LTR-3}{Early Warning Indicators}{STRIKE FIRST}{2}

\begin{lrkeybox}[Learning objectives — official 2026 spine wording]
\begin{enumerate}[label=\textbf{\alph*.}]
\item Evaluate the characteristics of sound Early Warning Indicators (EWI) measures.
\item Identify EWI guidelines from banking regulators and supervisors (OCC, BCBS, Federal Reserve).
\item Discuss the applications of EWIs in the context of the liquidity risk management process.
\end{enumerate}
{\footnotesize\color{FRMGrey}Source: Shyam Venkat and Stephen Baird, \textit{Liquidity Risk Management: A Practitioner's Perspective}, Chapter 6.}
\end{lrkeybox}

\begin{lrtrapbox}[Ordering note]
The Crash layer presents these objectives as a, c, b. They are reordered here into spine sequence. The Crash layer also runs objectives a and c under a single merged heading, which is why c looks thinner than it is in the source notes.
\end{lrtrapbox}

% =====================================================================
\lo{LTR-3 a}{Evaluate the characteristics of sound Early Warning Indicators (EWI) measures.}

One of the critical aspects of a bank's liquidity risk management (LRM) involves first devising and then monitoring a set of indicators, to enable the risk identification process to spot the emergence of new or increasing vulnerabilities.

\begin{lrkeybox}[Why EWIs matter]
\begin{itemize}
\item EWIs may be viewed as analogous to warning lights on an automobile dashboard. The turning signals, low oil, and check engine lights alert the driver of driving conditions.
\item Holding liquid assets to ensure a bank can meet its financial obligations is ultimately a \textbf{cost-benefit decision}.
\item Increasing buffers \emph{before} times of stress, with the help of EWIs, is cost-effective.
\end{itemize}
\end{lrkeybox}

\subsection{The EWI framework}

\gapnote{The source notes list the five components in this order but do not give them the mnemonic the GARP text uses. The acronym below is from the GARP curriculum; the five components and their content are the notes' own.}

\begin{lrdefbox}[M.E.R.I.T.]
The EWI framework can be summarized as \term{M.E.R.I.T.} — \textbf{M}easures, \textbf{E}scalation, \textbf{R}eporting, \textbf{I}ntegrated systems, and \textbf{T}hresholds. An appropriate set of measures is the first essential building block, but the framework is a mere academic exercise if it is not eventually linked to escalation processes. The journey from measures to escalation is facilitated by timely reporting, supported by integrated systems and data, and by relevant, properly calibrated thresholds.
\end{lrdefbox}

\begin{enumerate}
\item \textbf{Measures}: look at future liquidity risks, both on and off the balance sheet.
  \begin{enumerate}[label=\alph*)]
  \item Consider normal and stressed situations across different time horizons (hourly, daily, weekly, monthly) to forecast future cash flows and liquidity positions.
  \item EWIs should involve both \term{internal measures} (specific to the bank's balance sheet) and \term{external measures} (related to macroeconomic factors).
  \item EWIs should be \term{leading} (indicating problems before they occur) and \term{granular} (detailed enough to cut through noise).
  \end{enumerate}
\item \textbf{Escalation}: a clear plan to bring issues to relevant management based on severity.
\item \textbf{Reporting}: timely reports (daily, or even hourly for active trading) are crucial. Reports should be broad but focused on what matters most.
\item \textbf{Integrated systems}: an integrated data processing system allows more consistent and accurate reporting of measures when data is entered from multiple sources. Consistent and accurate data is key for good EWI reporting.
\item \textbf{Thresholds}: use a \term{green-amber-red} system (green = good, amber = investigate, red = urgent action).
  \begin{enumerate}
  \item Carefully set the amber zone to avoid missing problems or generating too many false alarms.
  \item Use past data to set thresholds and backtest them regularly for adjustments.
  \end{enumerate}
\end{enumerate}

\gapbadge
\gapnote{The notes state that EWIs should be leading and granular but do not define the two properties or name their source. Written from the GARP curriculum (Venkat and Baird, Chapter 6).}

\begin{lrgapbox}
The objective says \emph{evaluate} the characteristics. That requires knowing what makes an indicator leading rather than lagging, what \emph{sharpness} means with an example, and the BCBS principle the properties come from.
\end{lrgapbox}

\begin{center}
\small
\begin{tabularx}{\textwidth}{@{}p{2.6cm} X@{}}
\toprule
\rowcolor{FRMSoft}\textbf{Property} & \textbf{What it means, and the test} \\
\midrule
\textbf{Leading} & Provides information and signals potential stress \emph{prior to} the occurrence of an actual event — particularly important in preparing for systemic scenarios. A \textbf{lagging} indicator reports events that have already happened, such as government-reported GDP figures. A bank can instead build proxies for general economic performance from internal loan portfolio metrics. \\
\addlinespace
\textbf{Sharp} (granular) & The granularity and specificity of an indicator relative to the institution's profile. Sharp indicators are signals that do not go unnoticed within the mass of data. \textit{Example:} detecting a drop in \emph{overall} deposit balances is an acceptable EWI; detecting drops in the balances of \emph{more volatile segments} — high-net-worth customers, rate-sensitive products — brings into focus that certain important classes of customers are leaving the bank. \\
\addlinespace
\textbf{Internal vs external balance} & Internal measures are customized to the bank's balance sheet and activities; external measures signal systemic changes in the economy or market. Liquidity events can start \emph{either} inside the bank or be driven by the external environment — an idiosyncratic deposit run may result from the disclosure of poor performance \emph{or} from a systemic failure. \\
\bottomrule
\end{tabularx}
\end{center}
\figcap{Developing EWIs that are both sharp and forward-looking allows more time for management to take corrective or mitigating action, and gives managers a better understanding of risk drivers and trends than broad, lagging indicators provide.}

\begin{lrkeybox}[Where the properties come from: BCBS Principle 5]
Principle 5 of the BCBS \textit{Sound Principles} states that to obtain a forward-looking view of liquidity risk exposures, a bank should use metrics that assess \textbf{(a)} the structure of the balance sheet, as well as metrics that \textbf{(b)} project cash flows and future liquidity positions, taking into account \textbf{(c)} off-balance-sheet exposures.
\end{lrkeybox}

% =====================================================================
\lo{LTR-3 b}{Identify EWI guidelines from banking regulators and supervisors (OCC, BCBS, Federal Reserve).}

\gapbadge
\gapnote{The source notes reduce four regulatory documents to eight bullets in total. The table below is written from the GARP curriculum (Venkat and Baird, Chapter 6, Key Supervisory Guidelines). The notes' own eight items are all present within it, in bold.}

\begin{lrgapbox}
The objective says \emph{identify} the guidelines from three named supervisors. That means knowing which document says what, because a role-misassignment distractor here attributes one supervisor's guidance to another — and there are four documents, not three, since the BCBS appears twice.
\end{lrgapbox}

\begin{center}
\scriptsize
\begin{longtable}{@{}p{2.5cm} p{11.4cm}@{}}
\toprule
\rowcolor{FRMSoft}\textbf{Document} & \textbf{Major EWI components} \\
\midrule
\endfirsthead
\toprule
\rowcolor{FRMSoft}\textbf{Document} & \textbf{Major EWI components} \\
\midrule
\endhead
\textbf{OCC-2012} \newline {\color{FRMGrey}Liquidity booklet, OCC \textit{Comptroller's Handbook} (2012)} &
A bank should have EWIs that signal whether \textbf{embedded triggers} in certain products (for example callable public debt, OTC derivatives transactions) are about to be breached, or whether contingent risks are likely to materialize. \textbf{Early recognition of a potential event allows a bank to enhance its readiness.} EWIs may include:
\begin{itemize}
\item A reluctance of traditional fund providers to continue funding at historic levels
\item Pending regulatory action, formal or informal, or a CAMELS component or composite rating downgrade
\item Widening of spreads on senior and subordinated debt, credit default swaps, and stock price declines
\item Difficulty in accessing long-term debt markets
\item Reluctance of trust managers, money managers, public entities, and credit-sensitive fund providers to place funds
\item Rising funding costs in an otherwise-stable market
\item Counterparty resistance to off-balance-sheet products, or increased margin requirements
\item The elimination of committed credit lines by counterparties
\end{itemize} \\
\addlinespace
\textbf{BCBS-2008} \newline {\color{FRMGrey}``Principles for Sound Liquidity Risk Management and Supervision'' (2008)} &
A bank should design a set of indicators to identify the emergence of increased risk or vulnerabilities in its liquidity risk position or potential funding needs. EWIs can be \textbf{qualitative or quantitative} and may include, but are not limited to:
\begin{itemize}
\item \textbf{Rapid asset growth}, especially when funded with potentially volatile liabilities
\item \textbf{Growing concentrations in assets or liabilities}
\item Increases in currency mismatches
\item Decrease of weighted average maturity of liabilities
\item Repeated incidents of positions approaching or breaching internal or regulatory limits
\item Negative trends or heightened risk associated with a particular product line
\item \textbf{Significant deterioration in the bank's financial condition} (signalling deterioration of liquidity)
\item Negative publicity; credit rating downgrade; stock price declines; rising debt costs
\item Widening debt or credit-default-swap spreads; rising wholesale and retail funding costs
\item Counterparties requesting additional collateral or resisting new transactions; drop in credit lines
\item Increasing retail deposit outflows; increasing redemptions of CDs before maturity
\item Difficulty accessing longer-term funding; difficulty placing short-term liabilities
\item Increased concerns over the funding of off-balance-sheet items
\end{itemize} \\
\addlinespace
\textbf{BCBS-2012} \newline {\color{FRMGrey}``Monitoring Indicators for Intraday Liquidity'' (2012)} &
Intraday liquidity monitoring indicators include:
\begin{itemize}
\item \textbf{Daily maximum liquidity requirement}
\item \textbf{Available intraday liquidity}
\item Total payments
\item Time-specific and other critical obligations
\item Value of customer payments made on behalf of financial institution customers
\item Intraday credit lines extended to financial institution customers
\item Timing of intraday payments; intraday throughput
\end{itemize} \\
\addlinespace
\textbf{SR 10-6} \newline {\color{FRMGrey}Interagency Policy Statement on Funding and Liquidity Risk Management (2010), Federal Reserve Board} &
Institution management should monitor for potential liquidity stress events using early-warning indicators and \textbf{event triggers}, tailored to the institution's specific liquidity risk profile. Early recognition of potential events allows the institution to position itself into \textbf{progressive states of readiness} as the event evolves, while providing a framework to report or communicate within the institution and to outside parties. Early-warning signals may include, but are not limited to:
\begin{itemize}
\item \textbf{Negative publicity concerning an asset class owned by the institution}
\item \textbf{Increased potential for deterioration in the institution's financial condition}
\end{itemize} \\
\bottomrule
\end{longtable}
\end{center}
\figcap{Four documents, three supervisors — the BCBS appears twice, and its two appearances cover completely different ground. BCBS-2008 is the general firm-level EWI list; BCBS-2012 is intraday-only. A distractor that attributes ``daily maximum liquidity requirement'' to BCBS-2008, or ``rapid asset growth'' to BCBS-2012, is the exact role-misassignment shape this reading invites.}

\begin{lrtrapbox}[The discriminators worth memorising]
\begin{itemize}
\item \textbf{OCC-2012} is the one that talks about \textbf{embedded triggers in products} (callable debt, OTC derivatives).
\item \textbf{BCBS-2008} is the long general list, and the one that says EWIs may be \textbf{qualitative or quantitative}.
\item \textbf{BCBS-2012} is \textbf{intraday only}. If an option mentions daily maximum liquidity requirement or intraday throughput, it belongs here and nowhere else.
\item \textbf{SR 10-6} is the Federal Reserve one, and the one that introduces \textbf{event triggers} and \textbf{progressive states of readiness}.
\end{itemize}
\end{lrtrapbox}

% =====================================================================
\lo{LTR-3 c}{Discuss the applications of EWIs in the context of the liquidity risk management process.}

\gapbadge
\gapnote{The source notes cover this objective with a single line on industry practice. The placement of EWIs within the LRM process, below, is written from the GARP curriculum (Venkat and Baird, Chapter 6).}

\begin{lrgapbox}
The objective is about \emph{where EWIs sit} inside the wider liquidity risk management process — what feeds them, what they feed, and what regulation requires of that chain.
\end{lrgapbox}

It is important to ensure that EWIs align to, and are a natural extension of, the enterprise LRM framework. According to Federal Reserve Board Supervisory Letter \term{FRB SR 12-7}, the LRM framework should be end-to-end and ensure that the bank will sufficiently capture the banking organization's exposures, activities, and risks.

\begin{lrkeybox}[Where EWIs sit in the chain]
A comprehensive LRM framework runs in this order:
\begin{enumerate}
\item \textbf{Identify} the liquidity and funding risks inherent in the business activities of the institution.
\item \textbf{Memorialize} those risks in the firm's \term{liquidity risk inventory}.
\item \textbf{Assess} them through a risk assessment process involving the businesses, corporate treasury, the independent risk management function, and other relevant stakeholders.
\item \textbf{Measure and monitor} them through a framework consisting of \textbf{stress tests}, \textbf{EWIs}, and a \textbf{limit and response framework}.
\end{enumerate}
EWIs are therefore an \emph{output} of risk identification and assessment, not a substitute for it, and they sit alongside stress testing rather than replacing it.
\end{lrkeybox}

\subsection{Industry practices}

EWI dashboards are becoming more common, to improve risk reporting and supervisory duties.

\begin{lrtrapbox}[Consolidated trap summary — LTR-3]
\begin{itemize}
\item \textbf{Role.} Four documents, three supervisors. BCBS-2012 is intraday-only; BCBS-2008 is the general list. Attributing either one's content to the other is the primary distractor.
\item \textbf{Polarity.} \textbf{Leading} indicators signal before the event; \textbf{lagging} indicators report events that already happened. Government-reported GDP is the canonical lagging example.
\item \textbf{Scope.} \emph{Sharpness} is about granularity relative to the institution's own profile, not about statistical precision. Overall deposit balances are acceptable; volatile-segment balances are sharp.
\item \textbf{Role.} EWIs sit \emph{alongside} stress tests and the limit framework in the measure-and-monitor stage. They do not replace either, and they come after the risk inventory, not before it.
\item \textbf{Scope.} The amber threshold is the calibration problem — set too wide and problems are missed, too narrow and false alarms swamp the signal. Both failure directions matter.
\end{itemize}
\end{lrtrapbox}
